% GENERATED FILE. Edit canonical lessons, then run npm run generate:books.
% canonical-source-hash: fnv1a64:03cf32bb8a5ab878
% canonical-lessons: IT-C401-file, IT-C401-cancellare, IT-C401-copiare, IT-C401-incollare, IT-C401-backup

\chapter{File e cartelle}
\label{ch:it-file-e-cartelle}
\hlchaptermodality{401}

\begin{chapteropening}
\textbf{By the end of this chapter:} \emph{I can work with files: delete, copy and paste them, and make a backup.}

Complete the last lesson of chapter 401: I can work with files: delete, copy and paste them, and make a backup.
\end{chapteropening}

\section[il file]{il file — the computer file}

\label{lesson:IT-C401-file}



\begin{quote}
Before the new one: say the Italian for the device, then the Italian for the software.
\end{quote}

\subsection*{You'll want to know: il file}
\textbf{il file} — \textquotedblleft{}the computer file\textquotedblright{}.

\begin{quote}
Se mi mandi il file, lo stampo.
\end{quote}

\begin{quote}
\emph{If you send me the file, I will print it.}
\end{quote}

Masculine, invariable: \emph{i file}. Pronounced as in English.

\begin{cousinweb}
\emph{file} was borrowed from English \emph{file}, which came from French \emph{fil} \textquotedblleft{}thread\textquotedblright{}, from Latin \emph{fīlum} (Lewis \& Short): papers were once kept threaded on a string.
\end{cousinweb}

\begin{grammarlens}[title={What you've built}]
One more word for telling, explaining and describing at B1.
\end{grammarlens}

\subsection*{Your turn}
\begin{itemize}
\raggedright
  \item \emph{Say it:} \emph{il file}
  \item \emph{Say it:} \emph{il file}, once more
  \item \emph{Say it:} a sentence of your own with \emph{il file}
  \item \emph{Recall:} say the Italian for the device, then the Italian for the software, then say \emph{il file} again
\end{itemize}

\begin{tcolorbox}[breakable,colback=teal!4,colframe=teal!35!black,title={Before you move on}]
What does il file mean? (\textquotedblleft{}The computer file\textquotedblright{}.) Say it once more.
\end{tcolorbox}
\section[cancellare]{cancellare — to delete, to rub out; to cancel}

\label{lesson:IT-C401-cancellare}



\begin{quote}
Before the new one: say the Italian for the software, then the Italian for the computer file.
\end{quote}

\subsection*{You'll want to know: cancellare}
\textbf{cancellare} — \textquotedblleft{}to delete, to rub out; to cancel\textquotedblright{}.

\begin{quote}
Se cancelli la foto, non la trovi più.
\end{quote}

\begin{quote}
\emph{If you delete the photo, you will not find it again.}
\end{quote}

Also to cancel: \emph{il volo è stato cancellato}.

\begin{cousinweb}
\emph{cancellare} comes from Latin \emph{cancellāre} \textquotedblleft{}to make like a lattice, to cross out\textquotedblright{} (Lewis \& Short), from \emph{cancellī} \textquotedblleft{}a lattice, bars\textquotedblright{}: lines drawn across a text. The word \emph{cancello} \textquotedblleft{}gate\textquotedblright{} shares it.
\end{cousinweb}

\begin{grammarlens}[title={What you've built}]
One more word for telling, explaining and describing at B1.
\end{grammarlens}

\subsection*{Your turn}
\begin{itemize}
\raggedright
  \item \emph{Say it:} \emph{cancellare}
  \item \emph{Say it:} \emph{cancellare}, once more
  \item \emph{Say it:} a sentence of your own with \emph{cancellare}
  \item \emph{Recall:} say the Italian for the software, then the Italian for the computer file, then say \emph{cancellare} again
\end{itemize}

\begin{tcolorbox}[breakable,colback=teal!4,colframe=teal!35!black,title={Before you move on}]
What does cancellare mean? (\textquotedblleft{}To delete, to rub out; to cancel\textquotedblright{}.) Say it once more.
\end{tcolorbox}
\section[copiare]{copiare — to copy}

\label{lesson:IT-C401-copiare}



\begin{quote}
Before the new one: say the Italian for the computer file, then the Italian for to delete.
\end{quote}

\subsection*{You'll want to know: copiare}
\textbf{copiare} — \textquotedblleft{}to copy\textquotedblright{}.

\begin{quote}
Ho copiato il testo e te l'ho mandato.
\end{quote}

\begin{quote}
\emph{I copied the text and sent it to you.}
\end{quote}

At school, \emph{copiare} is to cheat in a test.

\begin{cousinweb}
\emph{copiare} is built on \emph{copia} \textquotedblleft{}copy\textquotedblright{}, from Medieval Latin \emph{cōpia}, which in Latin meant \textquotedblleft{}plenty, abundance\textquotedblright{} (Lewis \& Short): making many from one.
\end{cousinweb}

\begin{grammarlens}[title={What you've built}]
One more word for telling, explaining and describing at B1.
\end{grammarlens}

\subsection*{Your turn}
\begin{itemize}
\raggedright
  \item \emph{Say it:} \emph{copiare}
  \item \emph{Say it:} \emph{copiare}, once more
  \item \emph{Say it:} a sentence of your own with \emph{copiare}
  \item \emph{Recall:} say the Italian for the computer file, then the Italian for to delete, then say \emph{copiare} again
\end{itemize}

\begin{tcolorbox}[breakable,colback=teal!4,colframe=teal!35!black,title={Before you move on}]
What does copiare mean? (\textquotedblleft{}To copy\textquotedblright{}.) Say it once more.
\end{tcolorbox}
\section[incollare]{incollare — to paste; to glue}

\label{lesson:IT-C401-incollare}



\begin{quote}
Before the new one: say the Italian for to delete, then the Italian for to copy.
\end{quote}

\subsection*{You'll want to know: incollare}
\textbf{incollare} — \textquotedblleft{}to paste; to glue\textquotedblright{}.

\begin{quote}
Ho copiato l'indirizzo e l'ho incollato nel messaggio.
\end{quote}

\begin{quote}
\emph{I copied the address and pasted it into the message.}
\end{quote}

\emph{Copia e incolla} — copy and paste.

\begin{cousinweb}
\emph{incollare} is built from \emph{in-} and \emph{colla} \textquotedblleft{}glue\textquotedblright{}, which is inherited from Vulgar Latin \emph{colla} (EtymDB), from Greek \emph{kólla} \textquotedblleft{}glue\textquotedblright{}.
\end{cousinweb}

\begin{grammarlens}[title={What you've built}]
One more word for telling, explaining and describing at B1.
\end{grammarlens}

\subsection*{Your turn}
\begin{itemize}
\raggedright
  \item \emph{Say it:} \emph{incollare}
  \item \emph{Say it:} \emph{incollare}, once more
  \item \emph{Say it:} a sentence of your own with \emph{incollare}
  \item \emph{Recall:} say the Italian for to delete, then the Italian for to copy, then say \emph{incollare} again
\end{itemize}

\begin{tcolorbox}[breakable,colback=teal!4,colframe=teal!35!black,title={Before you move on}]
What does incollare mean? (\textquotedblleft{}To paste; to glue\textquotedblright{}.) Say it once more.
\end{tcolorbox}
\section[il backup]{il backup — the backup (copy)}

\label{lesson:IT-C401-backup}



\begin{quote}
Before the new one: say the Italian for to copy, then the Italian for to paste.
\end{quote}

\subsection*{You'll want to know: il backup}
\textbf{il backup} — \textquotedblleft{}the backup (copy)\textquotedblright{}.

\begin{quote}
Se fai un backup ogni settimana, non perdi niente.
\end{quote}

\begin{quote}
\emph{If you make a backup every week, you will not lose anything.}
\end{quote}

Masculine, invariable. Also \emph{la copia di sicurezza}.

\begin{cousinweb}
\emph{backup} was borrowed from English \emph{backup} (EtymDB), built from \emph{back} and \emph{up}: something kept in support.
\end{cousinweb}

\begin{grammarlens}[title={What you've built}]
One more word for telling, explaining and describing at B1.
\end{grammarlens}

\subsection*{Your turn}
\begin{itemize}
\raggedright
  \item \emph{Say it:} \emph{il backup}
  \item \emph{Say it:} \emph{il backup}, once more
  \item \emph{Say it:} a sentence of your own with \emph{il backup}
  \item \emph{Recall:} say the Italian for to copy, then the Italian for to paste, then say \emph{il backup} again
\end{itemize}

\begin{tcolorbox}[breakable,colback=teal!4,colframe=teal!35!black,title={Before you move on}]
What does il backup mean? (\textquotedblleft{}The backup (copy)\textquotedblright{}.) Say it once more.
\end{tcolorbox}
